\documentclass[11pt, a4paper]{article}
\usepackage[utf8]{inputenc}
\usepackage{amsmath, amssymb}
\usepackage[margin=1in]{geometry}
\usepackage{fancyhdr}
\usepackage{xcolor}
\definecolor{UNIBBlue}{RGB}{0, 51, 102}
\pagestyle{fancy}
\fancyhf{}
\rhead{\textbf{Optimisasi untuk Teknik Elektro}}
\lhead{Problem Set - Minggu 6}
\cfoot{\thepage}
\begin{document}
\begin{center}\Large\color{UNIBBlue}\textbf{Problem Set 6: Konveksitas}\end{center}
\section*{Soal 1 (C2)}
Jelaskan mengapa masalah optimisasi linier (LP) secara otomatis merupakan masalah optimisasi cembung!
\section*{Soal 2 (C4)}
Buktikan secara matematis apakah himpunan $S = \{ (x_1, x_2) \mid x_1^2 + x_2^2 \le 4 \}$ adalah himpunan cembung.
\section*{Soal 3 (C5)}
Diberikan sebuah masalah optimisasi untuk merancang antena. Anda menemukan bahwa fungsi biayanya non-cembung. Jelaskan risiko apa yang mungkin terjadi saat menggunakan algoritma \textit{Gradient Descent} biasa, dan apa solusinya?
\end{document}